\section{Introduction}
\label{sec:introduction}

Borsuk's partition problem asks whether every bounded set of positive
diameter in $\mathbb R^d$ can be covered by $d+1$ subsets of strictly
smaller diameter. For a bounded set $Y$ of positive diameter, write $b(Y)$ for the least
number of such subsets needed to cover $Y$. Covers and partitions give
the same number: overlaps may be removed without increasing diameters.
We prove that $b(Y)>10$ for a compact set in $\mathbb R^9$.

Let $\operatorname{Sym}_4(\mathbb R)$ be the space of real symmetric
$4\times4$ matrices, with Frobenius norm
$\|A\|_F^2=\operatorname{tr}(A^2)$. Its trace-one affine hyperplane is
a nine-dimensional Euclidean space. For a unit vector $u\in\mathbb R^4$,
the matrix $uu^{\mathsf T}$ is the orthogonal projector onto the line
spanned by $u$.

\begin{theorem}\label{thm:main}
The compact set
\[
 X=\{uu^{\mathsf T}:u\in\mathbb R^4,\ \|u\|=1\}
 \subset\{A\in\operatorname{Sym}_4(\mathbb R):\operatorname{tr}A=1\},
\]
with the Frobenius metric, has diameter $\sqrt2$ and cannot be covered
by ten subsets of diameter strictly less than $\sqrt2$.
\end{theorem}

Thus Borsuk's assertion fails in dimension nine. The theorem does not
determine the smallest dimension in which it fails
or the exact value of $b(X)$. Its witness is the full compact projector
image of real projective three-space. Corollary~\ref{cor:higher-dimensions}
extends the construction to compact counterexamples in every dimension
$d\geq9$.

\subsection*{History and the method}
Borsuk posed the partition question in 1933 \cite{Borsuk1933}.
Kahn and Kalai disproved it in high dimensions \cite{KahnKalai1993},
using a forbidden-intersection theorem of Frankl and Wilson
\cite{FranklWilson1981} to force many smaller-diameter parts.
Early reductions in counterexample dimension include $946$
by Nilli \cite{Nilli1994}, the announced bound $903$ of Grey and
Weissbach \cite{GreyWeissbach1997}, $561$ by Raigorodskii
\cite{Raigorodskii1997}, and $560$ by Weissbach \cite{Weissbach2000}.

Spherical codes led to a further series of reductions. Hinrichs obtained
dimension $323$ using vectors of the Leech lattice
\cite{Hinrichs2002}. Pikhurko reduced the bound to $321$ and credited
an independent discovery by Hinrichs and Richter \cite{Pikhurko2002};
Hinrichs and Richter published their bound $298$ in 2003
\cite{HinrichsRichter2003}.

A different series of advances used strongly regular graphs.
Bondarenko constructed a two-distance counterexample in dimension $65$
\cite{Bondarenko2014}; Jenrich and Brouwer obtained dimension $64$
\cite{JenrichBrouwer2014}. Grinsztajn's public proof note gives a
finite counterexample in dimension $63$ \cite{Grinsztajn2026}.
Theorem~\ref{thm:main} improves this dimension bound using the full
compact image of projective space rather than a finite configuration.

The projector embedding itself is classical. Kalai describes the map
$x\mapsto x\otimes x$ and the full rank-one positive semidefinite
configuration in his account of the Kahn--Kalai construction
\cite[Section~2.3]{Kalai2015}. Conway, Hardin, and Sloane use the
same projection-matrix model for Grassmannian packing and its Euclidean
metric \cite[Section~5, Theorem~2]{ConwayHardinSloane1996}.
Our contribution is the covering obstruction for the projector set
associated with $\mathbb R^4$.

The relevant topological antecedents include Walkup's eleven-vertex
minimum for triangulations of $\mathbb{RP}^3$
\cite[Theorem~3]{Walkup1970} and the projective-space vertex bounds
of Arnoux and Marin \cite{ArnouxMarin1991}. These results cannot be
applied directly to the complex here: its faces are positive coordinate
supports of a map, not simplices in a triangulation of projective space.
The central difficulty is to derive enough of its face structure from
the map itself. We do this through an odd map on symmetric matrices,
mod-two degree, and local preimage counts, and then give the finite
combinatorial obstruction in full.

\subsection*{Proof overview}
Two projectors in $X$ are at diameter distance precisely when their
underlying lines are orthogonal. A hypothetical cover by ten sets of
strictly smaller diameter can be enlarged to an open cover. A smooth
partition of unity then gives nonnegative functions $f_1,\ldots,f_{10}$
on $\mathbb{RP}^3$, with sum one, such that orthogonal lines have disjoint
sets of positive coordinates. Section~\ref{sec:geometry} proves this
reduction for arbitrary covering sets.

We record these sets of positive coordinates in a finite simplicial
complex: a set of labels is a face if all its coordinates are positive
at some line. To constrain this complex, Section~\ref{sec:extension}
extends the map $f$ to positive semidefinite matrices by integration,
then subtracts its values on the positive and negative parts of a
symmetric matrix. For the resulting odd map, the sum of the absolute output coordinates
equals the sum of the absolute eigenvalues of the input matrix.
Normalizing these two sums gives an odd map between spheres. In Section~\ref{sec:topology}, its mod-two
degree forces every maximal support to have four labels and the sum of
these tetrahedral faces to be a cycle.

The main topological step is the equality argument in
Section~\ref{sec:topology}. At the threshold of $k(k+1)/2$ labels on
$\mathbb{RP}^{k-1}$, a maximal support with fewer than $k$ labels
would give a positive-dimensional regular level set inside one affine
projective chart. Applying the matrix extension on the orthogonal
complement of each line gives a map on a sphere bundle over this level
set. The odd degree of the original extension forces an odd preimage
count for this map, whereas the chart trivialization makes the same
count even through the induced map on the projective bundle. Once every maximal support has
$k$ labels, the remaining odd local counts determine the coefficients
of the facet cycle.

The last stage converts this topological information into a finite
obstruction. The same degree argument in $\mathbb R^3$ produces a rigid triangle system on six
labels, described in Section~\ref{sec:six-labels}. Each tetrahedral face
of the ten-label complex forces such a system on its six-label
complement. Once these systems have been derived,
Section~\ref{sec:ten-labels} follows them across adjacent tetrahedra by
purely combinatorial arguments. Their compatibility restricts the cycles
around edges and organizes a subset of the labels into small blocks.
The blocks and the remaining labels give the final contradiction.
Section~\ref{sec:conclusion} returns to the diameter cover.
